% =====================================================================
%  Chapter 15 — Cosmological Probes
% =====================================================================

\chapter{Cosmological Probes and the Axis of Leverage}
\label{ch:cosmoprobes}

Cosmology has become the most powerful probe in the field, for a reason that
is not obvious from the outside. The other methods measure dark matter
through an interaction. Cosmology measures it through the effect of its
abundance on the growth of structure, which is a much larger effect and which
is sensitive to the mass of the particle rather than only to its coupling.

\section{The hierarchy of cosmological probes}Six probes matter, and they form a
hierarchy in what they constrain.

\begin{table}[H]
\centering
\caption{The cosmological probes, in decreasing order of the constraint they
impose on the dark matter particle mass. Each probes a different physical
effect, and the lower rows of the table are the ones that survive when the
upper rows are saturated.}
\label{tab:cosmoprobes}
\begin{tabular}{@{}>{\RaggedRight\arraybackslash}p{0.26\textwidth}p{0.31\textwidth}p{0.29\textwidth}@{}}
\toprule
\textbf{Probe} & \textbf{What it measures} & \textbf{Constraint} \\
\midrule
CMB acoustic peaks & $\Omega_c h^2$ and free streaming & $m > 0.3$--$1$~eV for thermal relics \\
Baryon acoustic oscillations & Distance scale and expansion history & Free-streaming scale at $z<2$ \\
Big bang nucleosynthesis & Baryon density at $z\sim 10^{9}$ & No dark matter effect \\
Linear large-scale structure & Amplitude and shape of the power spectrum & Momentum-dependent free streaming \\
Non-linear clustering and halo abundance & Small-scale power & $m > 10$--$100$~keV for thermal relics \\
Dwarf galaxies, satellites, cores & Structure at $10^{6}$--$10^{9}\ \Msun$ & Sub-GeV masses, warm and wave-like cases \\
\bottomrule
\end{tabular}
\end{table}

Two features of \tabref{tab:cosmoprobes} determine the strategy. The upper
probes are saturated: the CMB and BAO measurements of the dark matter
density are precise and will not improve much. The lower probes are where
the discriminating power now lives, and they are the ones that suffer most
from modelling uncertainty --- which is a statement about the data, not
about the physics.

\section{The free-streaming argument}The single most powerful cosmological
statement about dark matter is a negative one.

A particle with a mass $m$ and a velocity of order $v$ in the early universe
has a free-streaming length $\lambda_{\rm fs} \sim v/m$ (in natural units,
with a temperature-dependent factor). Free streaming erases structure below
that scale. The observed power spectrum has structure at scales corresponding
to a free-streaming length of about a kiloparsec, which for a thermal relic
corresponds to a mass of roughly 0.3 to 1~eV and rules out lighter thermal
relics entirely.

This is a much stronger constraint than the laboratory can offer for the same
mass range, because the cosmological effect is on the power spectrum itself
rather than on a rate of rare events. It also has a sharp consequence: the
thermal relic mass window is now narrow, with the interesting region being
near the electroweak scale and the interesting variants being the ones with
modified thermal histories.

\section{The mass function}The abundance of haloes of a given mass depends on the
shape of the power spectrum at small scales, and hence on the free-streaming
scale. This is what gives cosmology its reach into the keV and
sub-keV regimes, where laboratory experiments have essentially no
sensitivity and where the small-scale structure programme of
Chapter~\ref{ch:smallscale} has its discriminating power.

The difficulty is that halo abundances depend on the baryonic physics of
halo formation --- star formation, feedback, gas cooling --- and the
systematic uncertainty from that is comparable to or larger than the effect
being measured. This is the reason the halo abundance constraints on
warm dark matter have moved so little in a decade.

There are two escapes from this difficulty, and both are used. The first is
to use systems where baryonic physics is minimal: dwarf galaxies, where the
stellar feedback is weak, and clusters, where the baryon fraction can be
measured directly. The second is to use the shape of the concentration-mass
relation, which depends on the assembly history in a way that can be
modelled. Neither escape is complete, and the honest statement is that
cosmological constraints on the free-streaming scale have an uncertainty of
roughly a factor of two in mass, and that this factor is set by astrophysics
rather than by statistics.

\section{The new axis: weak lensing with spectroscopic redshifts}The most
significant development in the cosmological side of the field between 2024
and 2026 is not a new observable but a new combination. DESI's spectroscopic
redshifts, with more than 14 million galaxies over seven redshift bins, can
be combined with the wide-field imaging of the Dark Energy Survey and the
Hyper Suprime-Cam survey to measure the growth of structure with a
redshift resolution that photometric surveys alone cannot reach
\idx{DES Collaboration}.

This matters because it moves the cosmological tension described in
Chapter~\ref{ch:tensions} from the regime of a model-dependent comparison to
something closer to a direct measurement. If the growth of structure and the
expansion history are measured with the same galaxy population and the same
photometric systematics, then a discrepancy between them is much harder to
explain by calibration error.

\section{What cosmology cannot do}Cosmology measures the abundance of dark
matter, not its identity. A cosmological measurement can exclude a particle
mass, a coupling to baryons, or a production mechanism. It cannot detect a
particle, identify its spin, or determine its interaction with the Standard
Model. It also cannot distinguish a particle from a primordial black hole
population at any mass where the abundance is not otherwise constrained.

This is the structural reason cosmology and laboratory searches are
complements rather than substitutes. Chapter~\ref{ch:portfolio} turns this
into a programme.

\begin{keytakeaways}
\item Cosmological probes constrain the dark matter mass through its
free-streaming scale, which reaches far below the laboratory's reach but
with systematic uncertainties set by baryonic astrophysics.
\item The upper probes --- CMB and BAO --- are saturated; the discriminating
power has moved to the small-scale probes, where the modelling is worst.
\item Combining DESI spectroscopy with wide-field weak lensing converts a
tension between expansion and growth into something closer to a direct
measurement.
\item Cosmology constrains abundance and mass, not identity, which is why it
complements the laboratory rather than replacing it.
\end{keytakeaways}
